\appendix

\section{Source extraction for the LLTTY deterministic consequences}\label{app:LLTTY-extraction}

This appendix proves \cref{prop:LLTTY-package} as an author-derived result. Page and equation locators refer to \cite{LLTTY19}, arXiv:1801.05575v4. The table below separates exact source statements, immediate corollaries, and conclusions reconstructed from source proofs. No item is presented as a verbatim theorem unless it is one.

\small
\setlength{\LTleft}{0pt}
\setlength{\LTright}{0pt}
\setlength{\tabcolsep}{3pt}
\begin{longtable}{>{\raggedright\arraybackslash}p{0.13\linewidth}>{\raggedright\arraybackslash}p{0.24\linewidth}>{\raggedright\arraybackslash}p{0.20\linewidth}>{\raggedright\arraybackslash}p{0.35\linewidth}}
\toprule
\textbf{Item} & \textbf{Source locator} & \textbf{Status} & \textbf{Deduction recorded here}\\
\midrule
\endfirsthead
\toprule
\textbf{Item} & \textbf{Source locator} & \textbf{Status} & \textbf{Deduction recorded here}\\
\midrule
\endhead
\eqref{eq:class-cardinality} and class count
& pp.~34--35, Lemma~5.1, (14)--(16); p.~37, Lemmas~5.4--5.5
& direct restatement after reindexing
& Restriction to $\mathcal A_k(\mathcal X)$ cannot increase the number of nonempty equivalence classes.\\
$S_L=\bigcup K_u(L)\cup P_v(L)$
& p.~38, Theorem~5.6
& verbatim after notation change
& Codimension enters only through the previously defined gradual class.\\
\eqref{eq:factorial-compression}
& p.~55, Lemma~7.2
& stated ``in particular'' consequence
& Reindex the four source families by the manuscript's part index $q$.\\
\eqref{eq:small-bin-mass}
& p.~61, start of Case~2 in the proof of Lemma~7.7
& proof-level consequence
& Each of four geometrically decreasing families contributes $O(d^{-\kappa}n)$ below the cutoff.\\
\eqref{eq:K1} and \eqref{eq:K1-product}
& pp.~60--61, proof of Lemma~7.7, source (32)--(33), (38), Case~1
& author-derived packaging
& Multiply the three source families and use $n/s_q\le d^\kappa$.\\
\eqref{eq:K2} and \eqref{eq:K2-separation}
& p.~61, Case~2 in the proof of Lemma~7.7
& intermediate proof conclusion
& Define the first lower and last upper occupied weight orders and divide the total small-part mass by the per-order bound.\\
Low-order conclusion in part~(v)
& p.~62, Lemma~7.8
& direct source conclusion
& A repeated rounded value produces a weight set of mass at least $\delta/36$.\\
High-order conclusion and \eqref{eq:Pv-separation}
& pp.~62--63, Lemma~7.9 and its proof
& intermediate proof conclusion
& The height lower bound and $\wt w_q\ge h_q/d$ yield the separated upper order.\\
\bottomrule
\end{longtable}
\normalsize

Throughout the appendix, the source's gradual class $S$ is our codimension-dependent class $S_L$, and its families $K_u,P_u$ are denoted by $K_u(L),P_u(L)$. The codimension enters only through the definition of the steep classes and hence through $S_L$; the arguments in source Sections~5--7 are otherwise unchanged.

For a fixed canonical $k$-approximation, reindex all nonempty regular and spread $\ell$-parts as $L^{(q)}$, and write
\[
  s_q=\abs{L^{(q)}},
  \qquad
  h_q=h(L^{(q)}).
\]
The source's truncated weight $\wt w_q$ is exactly \eqref{eq:truncated-weights}; its $w$-set $W_b$ is the union of the parts with $2^b\le\wt w_q<2^{b+1}$. Thus the random dyadic order $B$ used here satisfies
\[
  \Pp\set{B=b}=\frac{\abs{W_b}}n.
\]
No rescaling or rebinning of the source weights is made.

\subsection{Equivalence classes and the gradual decomposition}

Source Lemma~5.1 states that, when $d\le n^{1/3}$ and
$1\le k\le\sqrt n/d^{3/2}$, the normalized canonical $k$-approximations form at most $e^n$ equivalence classes under equality of the coordinate multiset. This is precisely the relation $\sim_k$ on $\mathcal A_k(L)$. Restricting an equivalence relation to a subset cannot increase the number of nonempty classes, which gives the corresponding assertion for $\mathcal A_k(\mathcal X)$.

Source Lemma~5.4 shows that permutation of a canonical approximation preserves, order by order, the type, height, and cardinality of every nonempty $\ell$-part. Hence these data are well defined on each equivalence class. Source Lemma~5.5 then separates spread and regular parts by their orders and gives, in its notation,
\[
  \abs{\mathcal C}
  \le n!\prod_{j\ge0}
  \frac{(h_j^S)^{c_j^S}(h_j^R)^{c_j^R}}{c_j^S!c_j^R!}.
\]
After reindexing every nonempty spread or regular part by $q$, this is exactly
\[
  \abs{\mathcal C}
  \le\frac{n!}{\prod_qs_q!}\prod_qh_q^{s_q},
\]
which is \eqref{eq:class-cardinality}. Finally, source Theorem~5.6 states, for every integer $v\ge5$,
\[
  S_L=\bigcup_{u=4}^{v}K_u(L)\cup P_v(L).
\]
This proves parts~(i) and~(ii) of \cref{prop:LLTTY-package}.

\subsection{Factorial compression inside one weight order}

Fix a dyadic order $b$, put
\[
  I(b)=\set{q:2^b\le\wt w_q<2^{b+1}},
  \qquad
  t_b=\sum_{q\in I(b)}s_q=\abs{W_b}.
\]
Source Lemma~7.2 shows that the part sizes in a fixed $w$-set, arranged in nonincreasing order, are controlled by a geometric sequence. Its stated ``in particular'' conclusion is
\[
  \prod_{q\in I(b)}s_q!\ge t_b!e^{-Ct_b},
\]
which is \eqref{eq:factorial-compression}.

We also need the small-part consequence used in the second case of the proof of source Lemma~7.7. Let
\[
  T=d^{-\kappa}n,
  \qquad
  J(b)=\set{q\in I(b):s_q<T}.
\]
The proof of source Lemma~7.2 splits the parts in one $w$-set into four families according to regular/spread type and large/small status. Within each family, the sizes decrease geometrically with the order at which the parts were created. Consequently, the sum of the members of that family below the cutoff $T$ is $O(T)$. Summing the four families gives
\[
  \abs{\bigcup_{q\in J(b)}L^{(q)}}\le CT=Cd^{-\kappa}n.
\]
This is also the estimate written explicitly at the start of Case~2 in the proof of source Lemma~7.7. It proves part~(iii) of \cref{prop:LLTTY-package}.

\subsection{\texorpdfstring{The two $K_u$ alternatives}{The two Ku alternatives}}

Let $y$ be the $d^u$-approximation of a vector in $K_u(L)$. By the source definition of $K_u$, the union of its spread $\ell$-parts has size at least $\cK n_3$. Set
\[
  \kappa=\frac{\cK a_3}{16}<\frac13,
  \qquad
  \mathscr I=\set{q:s_q\ge d^{-\kappa}n}.
\]
Either
\[
  \sum_{q\in\mathscr I}s_q\ge n-\frac{\cK n_3}{4}
\]
or its complement has mass at least $\cK n_3/4$. These are alternatives \eqref{eq:K1} and \eqref{eq:K2}.

Assume \eqref{eq:K1}. Write $\mathscr I_1$ for the large spread parts, $\mathscr I_2$ for the large regular parts, and $\mathscr I_3=(\mathscr I_1\cup\mathscr I_2)^c$. Since the total spread mass is at least $\cK n_3$,
\[
  \sum_{q\in\mathscr I_1}s_q\ge\frac{3\cK n_3}{4},
  \qquad
  \sum_{q\in\mathscr I_3}s_q\le\frac{\cK n_3}{4}.
\]
For these three families, the definition of the truncated weights and the source lower bounds corresponding to its equations~(32)--(33) give
\begin{align*}
  \wt w_q^{-s_q}
  &\le d^{-s_q/2}\left(\frac n{s_q}\right)^{s_q/2}h_q^{-s_q},
  &&q\in\mathscr I_1,\\
  \wt w_q^{-s_q}
  &=\left(\frac n{s_q}\right)^{s_q}h_q^{-s_q},
  &&q\in\mathscr I_2,\\
  \wt w_q^{-s_q}
  &\le d^{s_q}h_q^{-s_q},
  &&q\in\mathscr I_3.
\end{align*}
Because $n/s_q\le d^\kappa$ on $\mathscr I_1\cup\mathscr I_2$, the same three-family calculation as in Case~1 of the proof of source Lemma~7.7 yields
\begin{align*}
  \prod_qh_q^{s_q}\prod_q\wt w_q^{-s_q}
  &\le C^n d^{-3\cK n_3/8}d^{\kappa n}d^{\cK n_3/4}\\
  &\le C^n d^{-\cK n_3/16}
  \le C^n d^{-c_1n}
\end{align*}
for an absolute $c_1>0$, after a harmless adjustment for $n_3=\floor{a_3n}$. This is \eqref{eq:K1-product}.

Assume now \eqref{eq:K2}. For each $b$, let
\[
  J(b)=\mathscr I^c\cap I(b).
\]
The preceding subsection gives
\[
  \abs{\bigcup_{q\in J(b)}L^{(q)}}\le C'd^{-\kappa}n.
\]
Following Case~2 in the proof of source Lemma~7.7, define $b_1$ to be the first order for which the cumulative lower $w$-sets have mass at least $\cK n_3/16$, and $b_2$ to be the last order for which the cumulative upper $w$-sets have the same mass. Then
\[
  \Pp\set{B\le b_1}\ge\frac{\cK n_3}{16n},
  \qquad
  \Pp\set{B\ge b_2}\ge\frac{\cK n_3}{16n}.
\]
The small parts between these two cutoffs have total mass at least $\cK n_3/8$. Since each intervening order contributes at most $C'd^{-\kappa}n$,
\[
  b_2-b_1+1
  \ge\frac{\cK n_3}{8C'd^{-\kappa}n}
  \ge cd^\kappa.
\]
After reducing fixed constants and taking $n$ large, this gives
\[
  \Pp\set{B\le b_1}\ge c_0,
  \qquad
  \Pp\set{B\ge b_2}\ge c_0,
  \qquad
  b_2-b_1\ge c_1d^\kappa.
\]
This proves part~(iv) of \cref{prop:LLTTY-package}.

\subsection{\texorpdfstring{The clustered $P_v$ alternative}{The clustered Pv alternative}}

Let $x\in P_v(L)$ have at least $\delta n$ coordinates in a disk of radius $\rho\le d^{-v}$, and let $y$ be its $d^v$-approximation. Source Lemma~7.8 shows that one value of $y$ is repeated on at least $\delta n/9$ coordinates and therefore produces a $w$-set of mass at least $\delta n/36$ and order at most
\[
  b_1=\log_2\left(\frac{72\sqrt d}{\delta}\right),
\]
with the integer rounding used in the source.

By the source definition of $P_v$, at least $\cP n_3$ coordinates belong to parts whose heights are at least
\[
  \cP a_3\,2^{\cP a_3(v-4)}.
\]
Since every truncated weight satisfies $\wt w_q\ge h_q/d$, these coordinates lie in $w$-sets above the high order used in the proof of source Lemma~7.9. The separation hypothesis \eqref{eq:Pv-separation-hyp} is exactly the hypothesis under which that proof establishes the intermediate conclusion
\[
  b_2-b_1\ge\frac{\cP a_3}{2}(v-4).
\]
Thus
\[
  \Pp\set{B\le b_1}\ge\frac\delta{36},
  \qquad
  \Pp\set{B\ge b_2}\ge\frac{\cP n_3}{n}\ge c_0
\]
for all sufficiently large $n$, after reducing the common absolute constant $c_0$ used above. The size condition \eqref{eq:Pv-delta} and the displayed separation condition are precisely the remaining assumptions of source Lemma~7.9. This proves part~(v) of \cref{prop:LLTTY-package} and completes the source extraction.
